% =====================================================================
% SECTION IV — METHODOLOGY
% File: section4_methodology.tex
% =====================================================================

% ------------------------------------------------------------------
\section{Proposed Methodology}
\label{sec:methodology}
% ------------------------------------------------------------------

We propose an Adaptive Temporal-Structural Fusion (ATSF) model
composed of five layers: (1) a Temporal Fusion Transformer branch,
(2) a Graph Attention Network branch, (3) a learned adaptive fusion
layer, (4) a demand output head, and (5) a rule-based risk scoring
interface. Fig.~\ref{fig:arch} illustrates the architecture.

% ------------------------------------------------------------------
\subsection{TFT Branch (Layer 1)}
% ------------------------------------------------------------------

The temporal branch processes historical demand sequences
$\mathbf{x} \in \mathbb{R}^{T \times d}$ together with static
covariates $\mathbf{s}$ (material type, normalised lead time). A
Variable Selection Network~\cite{lim2021tft} applies soft,
learned feature gating before a two-layer LSTM encoder (hidden
dimension $d{=}128$ for SupplyGraph, $d{=}64$ for USGS).
Multi-head self-attention ($H{=}4$ heads) over LSTM hidden states
captures long-range temporal dependencies.  Simultaneously, three
quantile heads ($q_{10}$, $q_{50}$, $q_{90}$) produce calibrated
prediction intervals via pinball loss. The branch outputs a
temporal embedding $\mathbf{h}_t \in \mathbb{R}^d$ and quantile
predictions $\hat{y}_{q} \in \mathbb{R}^{h}$, where $h$ is the
forecast horizon.

% ------------------------------------------------------------------
\subsection{GAT Branch (Layer 2)}
% ------------------------------------------------------------------

The structural branch encodes inter-commodity relationships through
a two-layer, multi-relational Graph Attention Network
\cite{velickovic2018gat}. Each node $v_i \in \mathcal{V}$ carries
four features: recent demand average, price-level proxy, production
utilisation, and normalised lead time. Edge attributes encode
relation type: 15 typed edges for USGS (e.g.\ \textit{supply\_chain\-input},
\textit{technology\_cluster}, \textit{battery\_coproduction}) and
4 typed edges for SupplyGraph (product group, sub-group, plant,
storage location). Typed attention gates prevent spurious
cross-relation aggregation~\cite{wasi2024gnn}. The branch outputs
a structural embedding $\mathbf{h}_s \in \mathbb{R}^d$.

% ------------------------------------------------------------------
\subsection{Adaptive Fusion Layer (Layer 3)}
% ------------------------------------------------------------------

The core contribution is a context-conditioned scalar fusion weight
$\alpha \in (0,1)$ learned end-to-end with the rest of the model.
Let $\mathbf{z}_m$ and $\mathbf{z}_h$ be 16-dimensional material-type
and horizon embeddings respectively. A residual LayerNorm projection
enriches the temporal embedding $\mathbf{h}_t$ into $\mathbf{h}_t'$;
this is concatenated with the structural embedding $\mathbf{h}_s$ and
both context embeddings and passed through a three-layer MLP to yield
$\alpha$ via a temperature-scaled sigmoid ($\tau{=}2.0$):

\begin{equation}
  \alpha = \sigma\!\Bigl(\mathrm{MLP}([\mathbf{h}_t',\,\mathbf{h}_s,\,
           \mathbf{z}_m,\,\mathbf{z}_h]) / \tau\Bigr)
  \label{eq:alpha}
\end{equation}

\begin{equation}
  \hat{\mathbf{h}} = \alpha\,\mathbf{h}_t + (1-\alpha)\,\mathbf{h}_s
  \label{eq:fused}
\end{equation}

\noindent The temperature $\tau{=}2.0$ softens the sigmoid, preventing
premature saturation during early training. A one-sided boundary
regularisation term

\begin{equation}
  \mathcal{L}_{\mathrm{reg}} = \lambda\!\left(
    \mathbb{E}[\mathrm{ReLU}(\alpha - 0.95)]
    + \mathbb{E}[\mathrm{ReLU}(0.05 - \alpha)]
  \right)
  \label{eq:reg}
\end{equation}

\noindent penalises collapse to either extreme whilst permitting
the model to discover a strong unimodal preference when the data
supports it ($\lambda{=}0.1$ SupplyGraph, $0.05$ USGS). This
design deliberately replaces entropy regularisation, which imposed
a centripetal bias toward $\alpha{=}0.5$. The fused representation
$\hat{\mathbf{h}}$ passes through a final linear-plus-LayerNorm
projection before the output head.

% ------------------------------------------------------------------
\subsection{Output and Risk Layers (Layers 4--5)}
% ------------------------------------------------------------------

A linear head maps $\hat{\mathbf{h}}$ to $h$ point forecast
values (Layer 4). Layer 5 applies four configurable, rule-based
thresholds — budget stress, lead-time risk, dependency criticality,
and a small learned risk calibrator MLP — to produce per-product
reorder flags from the forecast and graph attention outputs.

\smallskip
\noindent\textbf{Disclosure.} Layer 5 lead times for SupplyGraph
are synthetically generated (uniform random per product sub-group;
seed-controlled). Layer 5 thresholds are practitioner-configurable
and are not learned from data; this layer constitutes a
proof-of-concept decision-support interface and has not been
clinically validated.

\smallskip
\noindent\textbf{Training objective.} The total loss is
$\mathcal{L} = \mathcal{L}_{\mathrm{pinball}} +
               \mathcal{L}_{\mathrm{MSE}} +
               \mathcal{L}_{\mathrm{reg}}$,
where pinball loss across three quantiles serves the uncertainty
heads and MSE serves the point head. The model has approximately
198,000 trainable parameters.
